\documentclass[12pt,a4paper]{article}

\usepackage[utf8]{inputenc}
\usepackage{amsmath}
\usepackage{graphicx}
\usepackage{geometry}
\usepackage{hyperref}
\usepackage{booktabs}
\usepackage{caption}
\usepackage{float}
\usepackage{listings}
\usepackage{xcolor}
\usepackage{subcaption}

\geometry{margin=2.5cm}

\lstset{
    basicstyle=\footnotesize\ttfamily,
    breaklines=true,
    frame=single,
    numbers=left,
    numberstyle=\tiny,
    backgroundcolor=\color{gray!10},
}

\title{\textbf{Face Recognition Attendance System}\\[0.5cm]
\large{A Comparative Study of SVM and KNN Classifiers with FaceNet Embeddings}}
\author{
    Course: Image Processing and Recognition \\
    Project: Face Classification Using Cosine Similarity
}
\date{\today}

\begin{document}

\maketitle
\thispagestyle{empty}
\newpage

\tableofcontents
\newpage

\section{Problem Statement}
\label{sec:problem}

Traditional classroom attendance methods --- such as calling names, passing paper sign-in sheets, or using ID cards --- are time-consuming, prone to errors, and easily falsified. These manual approaches are inefficient for large classes and do not integrate naturally with modern digital workflows.

The goal of this project is to develop an automated, real-time face recognition attendance system with the following requirements:

\begin{enumerate}
    \item \textbf{Video-based enrollment}: Extract and enroll student faces from existing classroom videos, eliminating the need for dedicated enrollment sessions.
    \item \textbf{Data augmentation}: Synthesize additional training samples from limited video data to improve classifier robustness.
    \item \textbf{Classifier comparison}: Evaluate and compare two popular classification approaches --- Support Vector Machines (SVM) and K-Nearest Neighbors (KNN) --- on the same FaceNet embedding space to identify their respective strengths and weaknesses in a face recognition context.
    \item \textbf{Real-time inference}: Recognize students from a live webcam feed with minimal latency, using GPU acceleration.
    \item \textbf{Interactive interface}: Provide a user-friendly web interface for instructors to start/stop the camera, select the classifier, set confidence thresholds, and export attendance records.
\end{enumerate}

By comparing SVM and KNN on the same feature extraction backbone, this project aims to provide practical guidance on classifier selection for small-to-medium scale face recognition systems.

\section{Introduction}
\label{sec:intro}

Face recognition is one of the most important problems in computer vision, with applications ranging from security and surveillance to attendance tracking and human-computer interaction. This project aims to build a complete face recognition pipeline that extracts faces from classroom videos, trains classifiers on FaceNet embeddings, and deploys a real-time attendance system via a web interface.

The project is structured into four main phases:
\begin{enumerate}
    \item \textbf{Face Extraction}: Detect and crop faces from 22 classroom videos using MTCNN.
    \item \textbf{Data Augmentation}: Enrich the dataset with geometric and photometric transformations.
    \item \textbf{Model Training}: Extract 512-dimensional FaceNet embeddings and train two classifiers --- SVM (RBF kernel) and KNN (cosine distance) --- for comparison.
    \item \textbf{Real-time Application}: Deploy a Streamlit-based attendance system with live webcam feed, model selection, and CSV export.
\end{enumerate}

All deep learning inference (MTCNN, FaceNet) runs on GPU (CUDA 12.4) for real-time performance.

\subsection{Dataset}
\label{subsec:data}

The input dataset consists of 22 video files (\texttt{.mov}) of 22 students, stored in the \texttt{Video\_/} directory. Each video ranges from 5 to 25 seconds recorded at 30 fps. An additional 5 test videos (\texttt{.mp4}) are available in the \texttt{test-data/} directory.

The project name is \textbf{face-class-cosine}, and the codebase uses Python 3.10+ with PyTorch 2.6.0.

\newpage

\section{Face Extraction from Video}
\label{sec:extract}

Script: \texttt{extract\_faces.py}

\subsection{Methodology}
\label{subsec:extract_method}

Face extraction is performed using the following pipeline:
\begin{enumerate}
    \item Read each video from \texttt{Video\_/} using OpenCV.
    \item Sample frames at 1 fps (every 30 frames).
    \item Detect faces in each frame using MTCNN \cite{mtcnn}.
    \item Filter detections with confidence $> 0.95$.
    \item Crop and resize face regions to $160 \times 160$ pixels.
    \item Save to \texttt{dataset/train\_raw/\{student\_name\}/}.
\end{enumerate}

MTCNN (Multi-task Cascaded Convolutional Networks) is a state-of-the-art face detector that operates in three stages (P-Net, R-Net, O-Net) with high accuracy and GPU support. The image size of $160 \times 160$ matches the input requirement of FaceNet InceptionResNetV1.

\subsection{Results}
\label{subsec:extract_results}

A total of \textbf{320 face images} were extracted from 22 videos. Table~\ref{tab:extract_summary} shows the per-student distribution.

\begin{table}[H]
    \centering
    \caption{Face extraction statistics per student}
    \label{tab:extract_summary}
    \begin{tabular}{l r}
        \toprule
        \textbf{Student} & \textbf{Faces} \\
        \midrule
        daochitrung\_23ba14295      & 13 \\
        dohoanglong\_2410559        & 33 \\
        dohunganh\_23ba14011        & 24 \\
        duongtunganh\_2410024       & 10 \\
        kieuminhduc\_2411138        & 4  \\
        lehuuduc2411139             & 22 \\
        lethanhtra\_2410980         & 9  \\
        letrongtien\_23ba14280      & 8  \\
        luuminhhoang\_23BA14119     & 8  \\
        nguyenbinhduong\_22BA13091  & 13 \\
        nguyendinhgiang\_23ba14089  & 4  \\
        nguyendinhminhquang\_23BA14244 & 16 \\
        nguyenhong\_truong\_23BA14300 & 10 \\
        nguyenkhaiminh\_2410607     & 9  \\
        nguyenmy\_2410678           & 7  \\
        nguyenngochieu\_23BA14109   & 12 \\
        nguyennguyennhat\_23BA14222 & 44 \\
        nguyentathoangviet\_23BA14320 & 23 \\
        phanduyhoang\_23BA14117     & 17 \\
        phanminhtrang\_23BA14290    & 12 \\
        tranmanhhung\_23BA14127     & 10 \\
        transonbach\_2410136        & 12 \\
        \midrule
        \textbf{Total} & \textbf{320} \\
        \bottomrule
    \end{tabular}
\end{table}

Observation: The number of images per class is highly imbalanced (ranging from 4 to 44), which motivates the need for data augmentation.

\newpage

\section{Data Augmentation}
\label{sec:augment}

Script: \texttt{augment\_data.py}

\subsection{Methodology}
\label{subsec:augment_method}

To improve model generalization and balance the class distribution, we apply four transformations to each original image:

\begin{itemize}
    \item \textbf{Horizontal Flip}: Mirrors the image across the vertical axis.
    \item \textbf{Rotation}: Rotates by $\pm 10^\circ$ to simulate head tilts.
    \item \textbf{Brightness/Contrast Adjustment}: Alters brightness by factor $\alpha \in [0.8, 1.2]$ and contrast by $\beta \in [-30, 30]$.
    \item \textbf{Gaussian Blur}: Applies blur with kernel size 3 or 5 to simulate motion blur.
\end{itemize}

Each original image produces 4 augmented versions (one per transformation), increasing the dataset size by 5x.

\subsection{Results}
\label{subsec:augment_results}

After augmentation, the dataset grew from \textbf{320 to 1,600 images}. Figure~\ref{fig:dist_before} and Figure~\ref{fig:dist_after} show the distribution before and after augmentation.

\begin{figure}[H]
    \centering
    \includegraphics[width=\textwidth]{figures/distribution_before.pdf}
    \caption{Data distribution before augmentation (320 images, 22 classes)}
    \label{fig:dist_before}
\end{figure}

\begin{figure}[H]
    \centering
    \includegraphics[width=\textwidth]{figures/distribution_after.pdf}
    \caption{Data distribution after augmentation (1,600 images, 22 classes)}
    \label{fig:dist_after}
\end{figure}

The augmentation balanced the class sizes significantly (from 4--44 to 20--220 per class), enabling the model to learn better representations for under-represented classes.

\newpage

\section{Model Training}
\label{sec:training}

Two classifiers are trained independently in their respective directories:
\begin{itemize}
    \item \texttt{svm/train.py} --- FaceNet + SVM (RBF kernel)
    \item \texttt{knn/train.py} --- FaceNet + KNN (cosine distance)
\end{itemize}

\subsection{Common Pipeline}
\label{subsec:pipeline}

Both classifiers share a common feature extraction pipeline:

\begin{enumerate}
    \item \textbf{Feature Extraction}: FaceNet InceptionResNetV1 \cite{facenet}, pretrained on VGGFace2 \cite{vggface2}, extracts 512-dimensional embedding vectors for each face image. The model is run on GPU with \texttt{torch.no\_grad()} for efficiency.
    \item \textbf{Normalization}: Input images are normalized to $[-1, 1]$ using mean and standard deviation of 0.5.
    \item \textbf{Train/Test Split}: 80\% training, 20\% testing, stratified by class to preserve class proportions.
\end{enumerate}

FaceNet uses triplet loss to learn an embedding space where cosine distances between images of the same person are minimized and distances between different people are maximized. The resulting 512-dimensional vectors form a highly discriminative feature space.

\subsection{SVM Classifier}
\label{subsec:svm}

Script: \texttt{svm/train.py}

Support Vector Machine (SVM) is a powerful supervised learning algorithm for classification. The core idea is to find a hyperplane that separates classes with maximum margin.

In the 512-dimensional embedding space, data may not be linearly separable. We therefore use the \textbf{RBF (Radial Basis Function) kernel} to project data into a higher-dimensional feature space where classes become linearly separable. The RBF kernel is defined as:

\[
K(\mathbf{x}_i, \mathbf{x}_j) = \exp\left(-\gamma \|\mathbf{x}_i - \mathbf{x}_j\|^2\right)
\]

where $\gamma$ (gamma) controls the kernel width. A larger $\gamma$ means each training sample has a narrower influence radius.

The regularization parameter $C$ trades off between maximizing the margin and minimizing training errors:
\begin{itemize}
    \item Large $C$: prioritizes correct classification of all training samples, risking overfitting.
    \item Small $C$: allows misclassifications for a wider margin, improving generalization.
\end{itemize}

\textbf{Parameters}: $C = 10.0$, $\gamma = \text{scale}$ (automatically computed as $1 / (n\_features \times X.var())$), probability estimation enabled for confidence scoring.

\subsection{KNN Classifier}
\label{subsec:knn}

Script: \texttt{knn/train.py}

K-Nearest Neighbors (KNN) \cite{knn} is a lazy learning algorithm --- it stores all training data and defers computation until inference time. When predicting a new sample, KNN finds the $k$ closest training points in the feature space and votes to determine the label.

The distance metric is \textbf{cosine distance}:

\[
d(\mathbf{x}_i, \mathbf{x}_j) = 1 - \frac{\mathbf{x}_i \cdot \mathbf{x}_j}{\|\mathbf{x}_i\| \|\mathbf{x}_j\|}
\]

Cosine distance is well-suited for FaceNet embeddings because the embedding space is optimized for cosine similarity.

Voting uses \textbf{weights = distance}: closer neighbors have higher influence:

\[
P(y = c \mid \mathbf{x}) = \frac{\sum_{i \in \mathcal{N}_k(\mathbf{x})} w_i \cdot \mathbb{1}(y_i = c)}{\sum_{i \in \mathcal{N}_k(\mathbf{x})} w_i}
\quad\text{with}\quad w_i = 1 / d(\mathbf{x}, \mathbf{x}_i)
\]

The optimal $k$ is found via 5-fold cross-validation on the training set over $k \in \{1, 3, 5, 7, 9, 11, 15\}$.

\subsection{Results}
\label{subsec:training_results}

Table~\ref{tab:comparison} compares the two classifiers on the same test set (320 samples).

\begin{table}[H]
    \centering
    \caption{SVM vs KNN performance comparison}
    \label{tab:comparison}
    \begin{tabular}{l c c}
        \toprule
        \textbf{Metric} & \textbf{SVM (RBF)} & \textbf{KNN (k=1, cosine)} \\
        \midrule
        Test Accuracy (\%) & 99.69 & 99.69 \\
        Mean Confidence (\%) & 85.04 & 100.00 \\
        Training Time (s) & 2.38 & 0.00 \\
        Inference / 1000 samples (ms) & 0.77 & 0.21 \\
        Parameters & $C=10$, $\gamma=\text{scale}$ & $k=1$, cosine \\
        \bottomrule
    \end{tabular}
\end{table}

Both models achieve identical accuracy of \textbf{99.69\%} on the test set. Figure~\ref{fig:comparison} visualizes the differences across multiple metrics.

\begin{figure}[H]
    \centering
    \includegraphics[width=\textwidth]{figures/svm_vs_knn_comparison.pdf}
    \caption{SVM vs KNN comparison: accuracy/confidence (left), training/inference time (center), and KNN k-value cross-validation search (right).}
    \label{fig:comparison}
\end{figure}

\subsection{Strengths and Weaknesses Analysis}
\label{subsec:strengths_weaknesses}

Table~\ref{tab:swot} provides a comprehensive comparison of the two approaches.

\begin{table}[H]
    \centering
    \caption{Strengths and weaknesses of SVM and KNN}
    \label{tab:swot}
    \begin{tabular}{p{3.5cm} p{5.5cm} p{5.5cm}}
        \toprule
        \textbf{Criterion} & \textbf{SVM (RBF)} & \textbf{KNN (cosine)} \\
        \midrule
        \textbf{Strengths} &
        \begin{itemize}
            \item Handles non-linear data via kernel trick.
            \item Max-margin optimization improves generalization.
            \item Sparse decision space (support vectors) enables fast inference.
            \item Confidence scores reflect prediction uncertainty realistically.
        \end{itemize} &
        \begin{itemize}
            \item No training required (lazy learning) --- instant setup.
            \item Simple, interpretable, and easy to deploy.
            \item Naturally suited for well-structured embedding spaces.
            \item Trivially extensible with new data (no retraining).
            \item Faster inference than SVM for small class counts.
        \end{itemize} \\
        \midrule
        \textbf{Weaknesses} &
        \begin{itemize}
            \item Longer training time ($O(n^2)$ to $O(n^3)$).
            \item Hard to interpret: decisions depend on kernel and support vectors.
            \item Requires hyperparameter tuning ($C$, $\gamma$).
            \item Adding new data requires full retraining.
        \end{itemize} &
        \begin{itemize}
            \item Stores all training data (high memory usage).
            \item Slow inference on large datasets (brute-force neighbor search).
            \item Sensitive to noise and outliers with small $k$.
            \item 100\% confidence with $k=1$ does not reflect uncertainty.
            \item Performance degrades with high dimensionality (curse of dimensionality), though FaceNet embeddings mitigate this.
        \end{itemize} \\
        \bottomrule
    \end{tabular}
\end{table}

\newpage

\section{Embedding Space Analysis}
\label{sec:embedding}

Script: \texttt{analyze\_distribution.py}

To visualize the 512-dimensional embedding space, we apply PCA (Principal Component Analysis) to reduce dimensionality to 2 components. Figure~\ref{fig:embedding_pca} shows the distribution of embedding vectors on the PCA plane.

\begin{figure}[H]
    \centering
    \includegraphics[width=\textwidth]{figures/embedding_pca.pdf}
    \caption{PCA projection of face embedding vectors from test video}
    \label{fig:embedding_pca}
\end{figure}

The embeddings form well-separated clusters in 2D PCA space, confirming that FaceNet produces highly discriminative feature representations. Each cluster corresponds to a distinct individual, with minimal overlap between different people.

\newpage

\section{Real-time Attendance System}
\label{sec:streamlit}

Script: \texttt{app.py}

The real-time attendance system is built with Streamlit and offers the following features:

\begin{itemize}
    \item \textbf{Model Selection}: Radio button in the sidebar to choose between SVM (RBF) and KNN classifiers. Models are loaded dynamically from \texttt{svm/svm\_face\_classifier.pkl} and \texttt{knn/knn\_face\_classifier.pkl}.
    \item \textbf{Live Webcam Feed}: Captures video via OpenCV (DirectShow on Windows), processes frames at 5 fps with face detection downscaled to $320 \times 240$ pixels.
    \item \textbf{Face Detection and Recognition}: MTCNN detects faces; FaceNet extracts embeddings; the selected classifier predicts identity; boxes with confidence labels are drawn on the frame.
    \item \textbf{Attendance Tracking}: Students are marked present on their first detection with confidence $\ge$ threshold (default 0.5). Attendance is tracked per session using Streamlit's session state.
    \item \textbf{CSV Export}: Download attendance logs as CSV files with status, name, student ID, and timestamp.
    \item \textbf{GPU Acceleration}: All deep learning models run on CUDA. Heavy imports are loaded in a background thread to avoid UI lag, while \texttt{import torch} runs at module level for proper CUDA initialization.
\end{itemize}

\subsection{Architecture}
\label{subsec:architecture}

The application uses a singleton \texttt{CameraProcessor} class with three threads:
\begin{enumerate}
    \item \textbf{Capture thread}: Reads frames from the webcam at full framerate and stores the latest raw frame.
    \item \textbf{Processing thread}: Runs face detection, embedding extraction, and classification every $N=8$ frames. Results (JPEG bytes + attendance updates) are passed to the main thread.
    \item \textbf{Main (Streamlit) thread}: Renders the JPEG feed and attendance table using \texttt{st.image()} and \texttt{st.markdown()} with HTML tables (avoids the \texttt{st.dataframe} + \texttt{torch} crash bug on Windows).
\end{enumerate}

\newpage

\section{Project Structure}
\label{sec:project_structure}

The project is organized into a modular structure:

\begin{verbatim}
face-class-cosine/
├── svm/
│   ├── train.py                    # SVM training script
│   └── svm_face_classifier.pkl     # Trained SVM model
├── knn/
│   ├── train.py                    # KNN training script
│   └── knn_face_classifier.pkl     # Trained KNN model
├── extract_faces.py                # Face extraction from videos
├── augment_data.py                 # Data augmentation + distribution plots
├── analyze_distribution.py         # Regenerate all analysis figures
├── app.py                          # Streamlit real-time attendance app
├── train_comparison.py             # Comparative training (SVM + KNN)
├── dataset/
│   ├── train_raw/                  # Raw extracted faces (320 images)
│   └── train_augmented/            # Augmented dataset (1,600 images)
├── figures/                        # Generated figures (PDF)
│   ├── distribution_before.pdf
│   ├── distribution_after.pdf
│   ├── embedding_pca.pdf
│   └── svm_vs_knn_comparison.pdf
├── models/                         # Legacy model storage
├── report/
│   ├── report.tex                  # Vietnamese report
│   └── report_en.tex               # English report (this document)
├── test-data/                      # Test videos
└── Video_/                         # Original input videos
\end{verbatim}

\subsection{Script Descriptions}
\label{subsec:scripts}

\begin{table}[H]
    \centering
    \caption{Main scripts and their functions}
    \label{tab:scripts}
    \begin{tabular}{l l}
        \toprule
        \textbf{Script} & \textbf{Function} \\
        \midrule
        \texttt{extract\_faces.py} & Extract faces from videos using MTCNN \\
        \texttt{augment\_data.py}  & Data augmentation + distribution figures \\
        \texttt{svm/train.py}      & Train SVM classifier on FaceNet embeddings \\
        \texttt{knn/train.py}      & Train KNN classifier on FaceNet embeddings \\
        \texttt{train\_comparison.py} & Compare SVM and KNN side by side \\
        \texttt{analyze\_distribution.py} & Regenerate all distribution and PCA figures \\
        \texttt{app.py}           & Streamlit real-time attendance application \\
        \bottomrule
    \end{tabular}
\end{table}

\subsection{Libraries Used}
\label{subsec:libraries}

\begin{itemize}
    \item \textbf{PyTorch 2.6.0 + CUDA 12.4}: Deep learning framework for GPU-accelerated inference.
    \item \textbf{facenet-pytorch 2.5.3}: MTCNN face detection and FaceNet InceptionResNetV1 for embeddings.
    \item \textbf{OpenCV 4.13.0}: Video capture, image processing, and frame rendering.
    \item \textbf{scikit-learn 1.8.0}: SVM, KNN, PCA, cross-validation, and evaluation metrics.
    \item \textbf{Matplotlib 3.10.9 + Seaborn}: Data visualization and figure generation.
    \item \textbf{Streamlit}: Web application framework for the real-time attendance interface.
\end{itemize}

\section{Conclusion}
\label{sec:conclusion}

This project successfully built a complete face recognition pipeline from video data and compared two classifiers:

\begin{itemize}
    \item Extracted 320 face images from 22 classroom videos using MTCNN.
    \item Augmented the dataset to 1,600 images with four transformation types.
    \item Trained both SVM (RBF kernel, $C=10$) and KNN (cosine distance, $k=1$) on FaceNet embeddings, with both achieving \textbf{99.69\%} test accuracy.
    \item SVM provides more realistic confidence scores (85.04\%), while KNN offers instant training and faster inference.
    \item Deployed a real-time attendance system with Streamlit supporting both classifiers, model selection, and CSV export.
\end{itemize}

Future work includes exploring end-to-end models (ArcFace, CosFace) for potentially higher accuracy, integrating multi-view data, and optimizing KNN inference with KD-Tree or Ball Tree data structures for larger datasets.

\newpage

\begin{thebibliography}{9}
\bibitem{mtcnn} K. Zhang, Z. Zhang, Z. Li, and Y. Qiao, ``Joint Face Detection and Alignment using Multi-task Cascaded Convolutional Networks,'' \textit{IEEE Signal Processing Letters}, 2016.
\bibitem{facenet} F. Schroff, D. Kalenichenko, and J. Philbin, ``FaceNet: A Unified Embedding for Face Recognition and Clustering,'' \textit{CVPR}, 2015.
\bibitem{vggface2} Q. Cao, L. Shen, W. Xie, O. M. Parkhi, and A. Zisserman, ``VGGFace2: A dataset for recognising faces across pose and age,'' \textit{FG}, 2018.
\bibitem{knn} T. Cover and P. Hart, ``Nearest neighbor pattern classification,'' \textit{IEEE Transactions on Information Theory}, 1967.
\bibitem{svm} C. Cortes and V. Vapnik, ``Support-vector networks,'' \textit{Machine Learning}, 1995.
\end{thebibliography}

\end{document}
